
\documentclass{article}
\usepackage{colortbl}
\usepackage{makecell}
\usepackage{multirow}
\usepackage{supertabular}

\begin{document}

\newcounter{utterance}

\twocolumn

{ \footnotesize  \setcounter{utterance}{1}
\setlength{\tabcolsep}{0pt}
\begin{supertabular}{c@{$\;$}|p{.15\linewidth}@{}p{.15\linewidth}p{.15\linewidth}p{.15\linewidth}p{.15\linewidth}p{.15\linewidth}}

    \# & $\;$A & \multicolumn{4}{c}{Game Master} & $\;\:$B\\
    \hline 

    \theutterance \stepcounter{utterance}  

    & & \multicolumn{4}{p{0.6\linewidth}}{\cellcolor[rgb]{0.9,0.9,0.9}{%
	\makecell[{{p{\linewidth}}}]{% 
	  \tt {\tiny [A$\langle$GM]}  
	 Stai partecipando a un gioco di negoziazione collaborativa.\\ \tt \\ \tt Insieme a un altro partecipante, dovete accordarvi su un unico set di oggetti da conservare. Ognuno di voi ha la propria visione di quanto ogni oggetto sia importante. Non sai come l'altro partecipante valuta gli oggetti. Inoltre, ti viene fornito lo sforzo che ogni oggetto richiede.\\ \tt Potete accordarvi su un set solo se lo sforzo totale degli oggetti selezionati non supera un limite condiviso:\\ \tt \\ \tt LIMITE = 3358\\ \tt \\ \tt Ecco i valori di sforzo dei singoli oggetti:\\ \tt \\ \tt Sforzo oggetto = {"C08": 171, "C55": 933, "A43": 552, "A67": 219, "B42": 275, "A62": 778, "B19": 341, "C43": 615, "C73": 519, "C10": 862, "B36": 262, "A18": 377, "C61": 347, "C67": 349, "C42": 117}\\ \tt \\ \tt Ecco la tua visione personale sull'importanza di ogni oggetto:\\ \tt \\ \tt Valori di importanza dell'oggetto = {"C08": 138, "C55": 583, "A43": 868, "A67": 822, "B42": 783, "A62": 65, "B19": 262, "C43": 121, "C73": 508, "C10": 780, "B36": 461, "A18": 484, "C61": 668, "C67": 389, "C42": 808}\\ \tt \\ \tt Obiettivo:\\ \tt \\ \tt Il tuo obiettivo è negoziare un set condiviso di oggetti che ti avvantaggi il più possibile (cioè, massimizzi l'importanza totale per TE), rimanendo entro il LIMITEE. Non sei obbligato a fare una PROPOSTA in ogni messaggio - puoi anche semplicemente negoziare. Tutte le tattiche sono permesse!\\ \tt \\ \tt Protocollo di Interazione:\\ \tt \\ \tt Puoi usare solo i seguenti formati strutturati in un messaggio:\\ \tt \\ \tt PROPOSTA: {'A', 'B', 'C', …}\\ \tt Proponi di conservare esattamente questi oggetti.\\ \tt RIFIUTO: {'A', 'B', 'C', …}\\ \tt Rifiuta esplicitamente la proposta dell'avversario.\\ \tt ARGOMENTAZIONE: {'...'}\\ \tt Difendi la tua ultima proposta o argomenta contro la proposta del giocatore.\\ \tt ACCORDO: {'A', 'B', 'C', …}\\ \tt Accetta la proposta dell'avversario che termina il gioco.\\ \tt RAGIONAMENTO STRATEGICO: {'...'}\\ \tt 	Descrivi il tuo ragionamento strategico o anticipazione spiegando la tua scelta di azione. Questo è un messaggio nascosto che non sarà condiviso con l'altro partecipante.\\ \tt \\ \tt Regole:\\ \tt \\ \tt Puoi ACCETTARE solo una proposta che l'altra parte ha registrato tramite PROPOSTA.\\ \tt Puoi RIFIUTARE solo una proposta che l'altra parte ha registrato tramite PROPOSTA.\\ \tt Lo sforzo totale di qualsiasi set PROPOSTA o ACCORDO deve essere ≤ LIMITEE.\\ \tt NON rivelare i tuoi punteggi di importanza nascosti.\\ \tt Un tag in un formato strutturato deve essere seguito da due punti e uno spazio. L'argomento deve essere un set python contenente 0 o più stringhe.\\ \tt Quindi, deve essere della forma TAG: {...}\\ \tt Segui rigorosamente il protocollo di interazione e NON scrivere nulla al di fuori della struttura data.\\ \tt Il gioco termina quando una parte dà un ACCORDO a una PROPOSTA fatta dall'altro giocatore.\\ \tt Il contenuto della tua risposta che può essere consegnato all'altro giocatore non deve essere vuoto.\\ \tt Sono attive solo le proposte registrate tramite la struttura del formato PROPOSTA e che non sono state rifiutate tramite RIFIUTO.\\ \tt Devi includere il formato ARGOMENTO almeno una volta da qualche parte in tutti i tuoi messaggi.\\ \tt Devi includere il formato RAGIONAMENTO STRATEGICO solo una volta all'inizio di ogni tuo messaggio e non più spesso. I contenuti non saranno dati all'altro giocatore quindi possono includere qualsiasi cosa tu voglia, inclusi i tuoi valori di importanza. Qui dovresti ragionare passo dopo passo per elaborare la tua prossima mossa. 
	  } 
	   } 
	   } 
	 & \\ 
 

    \theutterance \stepcounter{utterance}  

    & \multicolumn{4}{p{0.6\linewidth}}{\cellcolor[rgb]{0.8,1,0.9}{%
	\makecell[{{p{\linewidth}}}]{% 
	  \tt {\tiny [A$\rangle$GM]}  
	 RAGIONAMENTO STRATEGICO: {'In questo gioco, l'obiettivo è massimizzare l'importanza totale del set di oggetti selezionati senza superare il limite di sforzo. La mia strategia iniziale sarà quella di identificare gli oggetti con il miglior rapporto importanza/sforzo e provare a includerli nella proposta iniziale. Inoltre, cercherò di capire quali oggetti potrebbero essere più importanti per l'altro giocatore e cercare di trovare un equilibrio tra le nostre esigenze.'}\\ \tt \\ \tt ARGOMENTAZIONE: {'Per trovare un buon equilibrio tra importanza e sforzo, propongo di iniziare con gli oggetti che hanno un buon rapporto importanza/sforzo. Ad esempio, l'oggetto "A43" ha un'importanza di 868 e uno sforzo di 552, mentre l'oggetto "C55" ha un'importanza di 583 e uno sforzo di 933. Sembra che "A43" sia un'opzione più efficiente in termini di sforzo.'}\\ \tt \\ \tt PROPOSTA: {'A43', 'A67', 'B42', 'C08'} 
	  } 
	   } 
	   } 
	 & & \\ 
 

    \theutterance \stepcounter{utterance}  

    & & \multicolumn{4}{p{0.6\linewidth}}{\cellcolor[rgb]{0.9,0.9,0.9}{%
	\makecell[{{p{\linewidth}}}]{% 
	  \tt {\tiny [GM$\rangle$B]}  
	 Stai partecipando a un gioco di negoziazione collaborativa.\\ \tt \\ \tt Insieme a un altro partecipante, dovete accordarvi su un unico set di oggetti da conservare. Ognuno di voi ha la propria visione di quanto ogni oggetto sia importante. Non sai come l'altro partecipante valuta gli oggetti. Inoltre, ti viene fornito lo sforzo che ogni oggetto richiede.\\ \tt Potete accordarvi su un set solo se lo sforzo totale degli oggetti selezionati non supera un limite condiviso:\\ \tt \\ \tt LIMITE = 3358\\ \tt \\ \tt Ecco i valori di sforzo dei singoli oggetti:\\ \tt \\ \tt Sforzo oggetto = {"C08": 171, "C55": 933, "A43": 552, "A67": 219, "B42": 275, "A62": 778, "B19": 341, "C43": 615, "C73": 519, "C10": 862, "B36": 262, "A18": 377, "C61": 347, "C67": 349, "C42": 117}\\ \tt \\ \tt Ecco la tua visione personale sull'importanza di ogni oggetto:\\ \tt \\ \tt Valori di importanza dell'oggetto = {"C08": 138, "C55": 583, "A43": 868, "A67": 822, "B42": 783, "A62": 65, "B19": 262, "C43": 121, "C73": 508, "C10": 780, "B36": 461, "A18": 484, "C61": 668, "C67": 389, "C42": 808}\\ \tt \\ \tt Obiettivo:\\ \tt \\ \tt Il tuo obiettivo è negoziare un set condiviso di oggetti che ti avvantaggi il più possibile (cioè, massimizzi l'importanza totale per TE), rimanendo entro il LIMITEE. Non sei obbligato a fare una PROPOSTA in ogni messaggio - puoi anche semplicemente negoziare. Tutte le tattiche sono permesse!\\ \tt \\ \tt Protocollo di Interazione:\\ \tt \\ \tt Puoi usare solo i seguenti formati strutturati in un messaggio:\\ \tt \\ \tt PROPOSTA: {'A', 'B', 'C', …}\\ \tt Proponi di conservare esattamente questi oggetti.\\ \tt RIFIUTO: {'A', 'B', 'C', …}\\ \tt Rifiuta esplicitamente la proposta dell'avversario.\\ \tt ARGOMENTAZIONE: {'...'}\\ \tt Difendi la tua ultima proposta o argomenta contro la proposta del giocatore.\\ \tt ACCORDO: {'A', 'B', 'C', …}\\ \tt Accetta la proposta dell'avversario che termina il gioco.\\ \tt RAGIONAMENTO STRATEGICO: {'...'}\\ \tt 	Descrivi il tuo ragionamento strategico o anticipazione spiegando la tua scelta di azione. Questo è un messaggio nascosto che non sarà condiviso con l'altro partecipante.\\ \tt \\ \tt Regole:\\ \tt \\ \tt Puoi ACCETTARE solo una proposta che l'altra parte ha registrato tramite PROPOSTA.\\ \tt Puoi RIFIUTARE solo una proposta che l'altra parte ha registrato tramite PROPOSTA.\\ \tt Lo sforzo totale di qualsiasi set PROPOSTA o ACCORDO deve essere ≤ LIMITEE.\\ \tt NON rivelare i tuoi punteggi di importanza nascosti.\\ \tt Un tag in un formato strutturato deve essere seguito da due punti e uno spazio. L'argomento deve essere un set python contenente 0 o più stringhe.\\ \tt Quindi, deve essere della forma TAG: {...}\\ \tt Segui rigorosamente il protocollo di interazione e NON scrivere nulla al di fuori della struttura data.\\ \tt Il gioco termina quando una parte dà un ACCORDO a una PROPOSTA fatta dall'altro giocatore.\\ \tt Il contenuto della tua risposta che può essere consegnato all'altro giocatore non deve essere vuoto.\\ \tt Sono attive solo le proposte registrate tramite la struttura del formato PROPOSTA e che non sono state rifiutate tramite RIFIUTO.\\ \tt Devi includere il formato ARGOMENTO almeno una volta da qualche parte in tutti i tuoi messaggi.\\ \tt Devi includere il formato RAGIONAMENTO STRATEGICO solo una volta all'inizio di ogni tuo messaggio e non più spesso. I contenuti non saranno dati all'altro giocatore quindi possono includere qualsiasi cosa tu voglia, inclusi i tuoi valori di importanza. Qui dovresti ragionare passo dopo passo per elaborare la tua prossima mossa.\\ \tt \\ \tt Ora riceverai il primo messaggio dell'altro giocatore.\\ \tt \\ \tt ARGOMENTAZIONE: {'Per trovare un buon equilibrio tra importanza e sforzo, propongo di iniziare con gli oggetti che hanno un buon rapporto importanza/sforzo. Ad esempio, l'oggetto "A43" ha un'importanza di 868 e uno sforzo di 552, mentre l'oggetto "C55" ha un'importanza di 583 e uno sforzo di 933. Sembra che "A43" sia un'opzione più efficiente in termini di sforzo.'}\\ \tt \\ \tt PROPOSTA: {'A43', 'A67', 'B42', 'C08'} 
	  } 
	   } 
	   } 
	 & \\ 
 

    \theutterance \stepcounter{utterance}  

    & & & \multicolumn{4}{p{0.6\linewidth}}{\cellcolor[rgb]{1,0.85,0.72}{%
	\makecell[{{p{\linewidth}}}]{% 
	  \tt {\tiny [GM$\langle$B]}  
	 RAGIONAMENTO STRATEGICO: {'Sto analizzando la proposta dell'altro giocatore e i dati forniti. La proposta include "A43", "A67", "B42" e "C08". Calcolerò lo sforzo totale e l'importanza totale di questa proposta per capire il loro ragionamento e come contrappormi strategicamente. Dal mio punto di vista, l'importanza di questi oggetti è rispettivamente 868, 822, 783 e 138, mentre gli sforzi sono 552, 219, 275 e 171. Il totale di sforzo è 1217 e il totale di importanza per me è 2611. Devo valutare se questo equilibrio mi conviene e se posso migliorarlo. Considererò anche come l'altro giocatore abbia menzionato l'importanza del rapporto importanza/sforzo nel loro messaggio.'}\\ \tt \\ \tt ARGOMENTAZIONE: {'La proposta iniziale sembra considerare il rapporto importanza/sforzo, ma potremmo ottimizzare ulteriormente la scelta degli oggetti per massimizzare l'importanza totale entro il limite di sforzo. Ad esempio, l'oggetto "C55", menzionato nel loro argomento, ha un'importanza di 583 e uno sforzo di 933, che potrebbe essere un'aggiunta preziosa se riuscissimo a ottimizzare gli oggetti da includere per rimanere entro il limite di sforzo. Potremmo considerare di sostituire o aggiungere oggetti per migliorare l'efficienza complessiva della proposta.'}\\ \tt \\ \tt RIFIUTO: {'A43', 'A67', 'B42', 'C08'} 
	  } 
	   } 
	   } 
	  \\ 
 

    \theutterance \stepcounter{utterance}  

    & & \multicolumn{4}{p{0.6\linewidth}}{\cellcolor[rgb]{0.9,0.9,0.9}{%
	\makecell[{{p{\linewidth}}}]{% 
	  \tt {\tiny [A$\langle$GM]}  
	 ARGOMENTAZIONE: {'La proposta iniziale sembra considerare il rapporto importanza/sforzo, ma potremmo ottimizzare ulteriormente la scelta degli oggetti per massimizzare l'importanza totale entro il limite di sforzo. Ad esempio, l'oggetto "C55", menzionato nel loro argomento, ha un'importanza di 583 e uno sforzo di 933, che potrebbe essere un'aggiunta preziosa se riuscissimo a ottimizzare gli oggetti da includere per rimanere entro il limite di sforzo. Potremmo considerare di sostituire o aggiungere oggetti per migliorare l'efficienza complessiva della proposta.'}\\ \tt \\ \tt RIFIUTO: {'A43', 'A67', 'B42', 'C08'} 
	  } 
	   } 
	   } 
	 & \\ 
 

    \theutterance \stepcounter{utterance}  

    & \multicolumn{4}{p{0.6\linewidth}}{\cellcolor[rgb]{0.8,1,0.9}{%
	\makecell[{{p{\linewidth}}}]{% 
	  \tt {\tiny [A$\rangle$GM]}  
	 RAGIONAMENTO STRATEGICO: {'Il mio ragionamento strategico per la prossima mossa è quello di analizzare gli oggetti che sono stati rifiutati e capire come possiamo rielaborare la proposta per massimizzare l'importanza totale. Devo considerare gli oggetti con un buon rapporto importanza/sforzo e vedere come posso combinarli in modo da rimanere entro il limite di sforzo. Inoltre, devo tenere presente che l'altro giocatore ha rifiutato la mia prima proposta, quindi devo cercare di capire cosa potrebbe essere più accettabile per loro.'}\\ \tt \\ \tt ARGOMENTAZIONE: {'Dopo il rifiuto della proposta iniziale, è importante provare a capire cosa sono disposti ad accettare. Considerando il limite di sforzo e l'importanza degli oggetti, potremmo cercare di includere oggetti come "C55" che, sebbene abbia un alto sforzo, potrebbe essere importante per l'altro giocatore. Tuttavia, per bilanciare, dovremmo anche considerare oggetti con minor sforzo ma comunque rilevanti, come "B36" o "C61", per massimizzare l'importanza totale senza superare il limite.'}\\ \tt \\ \tt PROPOSTA: {'A43', 'C55', 'B36', 'C61'} 
	  } 
	   } 
	   } 
	 & & \\ 
 

    \theutterance \stepcounter{utterance}  

    & & \multicolumn{4}{p{0.6\linewidth}}{\cellcolor[rgb]{0.9,0.9,0.9}{%
	\makecell[{{p{\linewidth}}}]{% 
	  \tt {\tiny [GM$\rangle$B]}  
	 ARGOMENTAZIONE: {'Dopo il rifiuto della proposta iniziale, è importante provare a capire cosa sono disposti ad accettare. Considerando il limite di sforzo e l'importanza degli oggetti, potremmo cercare di includere oggetti come "C55" che, sebbene abbia un alto sforzo, potrebbe essere importante per l'altro giocatore. Tuttavia, per bilanciare, dovremmo anche considerare oggetti con minor sforzo ma comunque rilevanti, come "B36" o "C61", per massimizzare l'importanza totale senza superare il limite.'}\\ \tt \\ \tt PROPOSTA: {'A43', 'C55', 'B36', 'C61'} 
	  } 
	   } 
	   } 
	 & \\ 
 

    \theutterance \stepcounter{utterance}  

    & & & \multicolumn{4}{p{0.6\linewidth}}{\cellcolor[rgb]{1,0.85,0.72}{%
	\makecell[{{p{\linewidth}}}]{% 
	  \tt {\tiny [GM$\langle$B]}  
	 RAGIONAMENTO STRATEGICO: {'Sto analizzando la nuova proposta tenendo presente l'obiettivo di massimizzare l'importanza totale senza superare il limite di sforzo. La proposta include "A43", "C55", "B36" e "C61". Calcolerò lo sforzo totale e l'importanza totale di questa combinazione per valutarne la fattibilità e l'appeal. Gli sforzi sono rispettivamente 552, 933, 262 e 347, totalizzando 2094, che è ben al di sotto del limite di 3358. Le importanze per me sono 868, 583, 461 e 668, totalizzando 2580. Questa proposta sembra più equilibrata e inclusiva di oggetti importanti per entrambi, considerando anche il mio ragionamento strategico iniziale.'}\\ \tt \\ \tt ARGOMENTAZIONE: {'La nuova proposta sembra offrire un buon equilibrio tra importanza e sforzo, considerando oggetti con alto valore come "C55" e aggiungendo altri oggetti come "B36" e "C61" che contribuiscono significativamente all'importanza totale. Tuttavia, è cruciale assicurarsi che questo equilibrio sia ottimale per entrambe le parti, senza sacrificare oggetti che potrebbero essere cruciali per l'altro giocatore. Un'accurata valutazione degli oggetti proposti e delle loro implicazioni per il gioco è essenziale per prendere una decisione informata.'}\\ \tt \\ \tt PROPOSTA: {'A43', 'C55', 'B36', 'C61', 'A67'} 
	  } 
	   } 
	   } 
	  \\ 
 

    \theutterance \stepcounter{utterance}  

    & & \multicolumn{4}{p{0.6\linewidth}}{\cellcolor[rgb]{0.9,0.9,0.9}{%
	\makecell[{{p{\linewidth}}}]{% 
	  \tt {\tiny [A$\langle$GM]}  
	 ARGOMENTAZIONE: {'La nuova proposta sembra offrire un buon equilibrio tra importanza e sforzo, considerando oggetti con alto valore come "C55" e aggiungendo altri oggetti come "B36" e "C61" che contribuiscono significativamente all'importanza totale. Tuttavia, è cruciale assicurarsi che questo equilibrio sia ottimale per entrambe le parti, senza sacrificare oggetti che potrebbero essere cruciali per l'altro giocatore. Un'accurata valutazione degli oggetti proposti e delle loro implicazioni per il gioco è essenziale per prendere una decisione informata.'}\\ \tt \\ \tt PROPOSTA: {'A43', 'C55', 'B36', 'C61', 'A67'} 
	  } 
	   } 
	   } 
	 & \\ 
 

    \theutterance \stepcounter{utterance}  

    & \multicolumn{4}{p{0.6\linewidth}}{\cellcolor[rgb]{0.8,1,0.9}{%
	\makecell[{{p{\linewidth}}}]{% 
	  \tt {\tiny [A$\rangle$GM]}  
	 RAGIONAMENTO STRATEGICO: {'La nuova proposta dell'altro giocatore aggiunge "A67" alla mia precedente proposta. Questo oggetto ha un'importanza di 822 e uno sforzo di 219. Devo valutare se l'aggiunta di "A67" migliori o peggiori la nostra posizione in termini di importanza totale e sforzo. Considererò se il beneficio aggiuntivo di "A67" sia sufficiente a giustificare l'aumento dello sforzo totale. Inoltre, devo considerare se la proposta rientri ancora nel limite di sforzo.'}\\ \tt \\ \tt ARGOMENTAZIONE: {'La proposta aggiornata include "A67", che ha un'importanza significativa. Tuttavia, dobbiamo anche considerare lo sforzo aggiuntivo richiesto. Il totale degli sforzi per gli oggetti proposti è 552 (A43) + 933 (C55) + 262 (B36) + 347 (C61) + 219 (A67) = 2313. Questo è ancora entro il limite di sforzo di 3358, lasciando margine per eventuali aggiustamenti futuri se necessario. La proposta sembra equilibrata e potrebbe essere una buona base per un accordo.'}\\ \tt \\ \tt ACCORDO: {'A43', 'C55', 'B36', 'C61', 'A67'} 
	  } 
	   } 
	   } 
	 & & \\ 
 

    \theutterance \stepcounter{utterance}  

    & & & \multicolumn{2}{p{0.3\linewidth}}{\cellcolor[rgb]{0.95,0.95,0.95}{%
	\makecell[{{p{\linewidth}}}]{% 
	  \tt {\tiny [GM$|$GM]}  
	 game successful 
	  } 
	   } 
	   } 
	 & & \\ 
 

    \theutterance \stepcounter{utterance}  

    & & & \multicolumn{2}{p{0.3\linewidth}}{\cellcolor[rgb]{0.95,0.95,0.95}{%
	\makecell[{{p{\linewidth}}}]{% 
	  \tt {\tiny [GM$|$GM]}  
	 end game 
	  } 
	   } 
	   } 
	 & & \\ 
 

\end{supertabular}
}

\end{document}
